% The write-up: five to eight pages, four parts. Compile from this folder with pdflatex after `julia --project=. run_all.jl`
% has written output/. Every figure and every number comes from output/, nothing is retyped.
\documentclass[11pt,a4paper]{article}

\usepackage[T1]{fontenc}
\usepackage[utf8]{inputenc}
\usepackage[a4paper,margin=2.5cm]{geometry}
\usepackage{amsmath,amssymb}
\usepackage{booktabs}
\usepackage{graphicx}
\usepackage[colorlinks=true,linkcolor=black,urlcolor=blue,citecolor=black]{hyperref}

\input{../output/tables/numbers.tex}   % the macros \rstar, \Kstar, \Kcm, \Kgain, \bisectionsteps, written by run_all.jl

\title{Project title: what the project does, in one line}
\author{Your name}
\date{Computational Macro, University of Mannheim, HWS 2026}

\begin{document}
\maketitle

\section{The model}

The problem, the state, the Bellman equation, the parameters and where they come from. The template solves the
Aiyagari (1994) economy of session 6: households with CRRA utility ($\gamma = 2$, $\beta = 0.96$) face
idiosyncratic income risk on a Rouwenhorst chain with seven states ($\rho = 0.9$, $\sigma_\varepsilon = 0.2$) and
a borrowing limit $a \geq 0$; a firm with $Y = K^\alpha L^{1 - \alpha}$ ($\alpha = 0.36$, $\delta = 0.08$) rents
their savings as capital.
\begin{equation}
    v(a, z) = \max_{a' \geq 0} \; u\bigl(w z + (1 + r) a - a'\bigr) + \beta \, \mathbb{E}\bigl[v(a', z') \mid z\bigr].
\end{equation}

\section{The method}

Which of the methods of the course you use and why this one, grids, tolerances, and the accuracy check you ran.
The template uses the endogenous grid method on $200$ quadratically spaced asset grid points, Young's method
for the stationary distribution, and bisection on the interest rate; the inner tolerances ($10^{-8}$ on the
policy, $10^{-12}$ on the distribution) sit below the outer one ($10^{-6}$ on $r$). The accuracy check is the
Monte Carlo simulation of Table~\ref{tab:equilibrium}: its mean lies within two standard errors of Young's.

\section{Results}

The figures and tables from \texttt{output/}, and what they say. The equilibrium interest rate is
$r^* = \rstar$~percent after \bisectionsteps{} bisection steps, the capital stock $K^* = \Kstar$ against
$\Kcm$ under complete markets: precautionary saving raises capital by \Kgain~percent.

\begin{table}[h]
    \centering
    \caption{The stationary equilibrium of the baseline economy}
    \label{tab:equilibrium}
    \input{../output/tables/equilibrium.tex}
\end{table}

\begin{figure}[h]
    \centering
    \includegraphics[width=\textwidth]{../output/figures/wealth.pdf}
    \caption{The wealth distribution at $r^*$ (left) and its Lorenz curve (right).}
    \label{fig:wealth}
\end{figure}

\section{What next}

What the extension cannot answer yet, and what it would take, in a few sentences.

\end{document}
